% ── 3. SenseMath Benchmark ────────────────────────────────────────────────────
\section{The \textsc{SenseMath} Benchmark}
\label{sec:benchmark}

\textsc{SenseMath} is a controlled benchmark for measuring whether LLMs can exploit number-sense shortcuts. This section describes our design philosophy, category taxonomy, programmatic generation pipeline, and evaluation framework.

\subsection{Design Philosophy}
\label{sec:design}

The central question is whether LLMs can reduce computation when a problem's numerical structure permits it. We operationalise this through three design choices.

\textbf{Shortcut definition.}
A shortcut is \emph{effective} if it reduces the problem to operations trivially executable in one's head: single-digit arithmetic, multiplication by powers of 10, or simple addition after simplification.
A shortcut is \emph{not} effective if the simplified form still requires multi-digit multiplication (e.g., rounding $5{,}374 \times 2{,}169$ to $5{,}400 \times 2{,}200$ does not meaningfully reduce computation), or if the approximation cannot uniquely identify the correct answer among the given options.

\textbf{Matched family design.}
Each item belongs to a \emph{family} containing three variants: \emph{strong} (clean shortcut available), \emph{weak} (partial shortcut requiring more effort), and \emph{control} (no effective shortcut). All three share the same template and digit scale, differing only in numerical values. This matched design lets us attribute accuracy differences to shortcut availability rather than problem difficulty.

\textbf{Three-level evaluation.}
We evaluate number sense at three levels of cognitive demand:
\emph{Use} (can the model solve the problem more accurately when encouraged to use shortcuts?),
\emph{Judge} (can the model recognise when a shortcut applies or identify a shortcut strategy in a solution?), and
\emph{Generate} (can the model construct new shortcut/control problem pairs?).

\subsection{Categories}
\label{sec:categories}

\textsc{SenseMath} spans 8 categories organised into three tiers based on where the shortcut operates.

\textbf{Tier 1: Problem-level shortcuts} target competencies from the mathematics education literature \citep{mcintosh1992proposed,reys1991number}.
\emph{Magnitude estimation (ME)}: rounding operands to nearby powers of 10.
\emph{Structural shortcuts (SS)}: exploiting algebraic identities near round numbers (e.g., $99 \times 37 = (100{-}1) \times 37$).
\emph{Relative distance (RD)}: comparing fractions by their distance from a benchmark (e.g., $\nicefrac{1}{2}$).
\emph{Cancellation (CI)}: near-cancellation in $A + B - C$ when $B \approx C$.
\emph{Compatible numbers (CN)}: rounding to product-friendly pairs (e.g., $250 \times 4000$).
\emph{Landmark comparison (LC)}: comparing percentages to landmarks (e.g., $49\% \approx 50\%$).

\textbf{Tier 2: Reasoning-level shortcut.}
\emph{Equation reasoning (ER)}: recognising algebraic identities (commutativity, common-term cancellation) in fill-in-the-blank equations, reducing multi-step algebra to a single observation.

\textbf{Tier 3: Option-level shortcut.}
\emph{Option elimination (OE)}: ruling out answer choices by quick feasibility checks (trailing digit, parity, magnitude) without computing the exact answer.

Example items for each category are provided in Appendix~\ref{sec:appendix_examples}.

\subsection{Program-Based Generation}
\label{sec:generation}

All items are generated programmatically via category-specific Python generators with rejection sampling. This ensures:
(i) 100\% answer correctness by construction (exact arithmetic at generation time);
(ii) distractors that share the last 50\% of digits with the correct answer (arithmetic categories) or fall within $\pm 0.1$ decimal (fraction categories), preventing option-level elimination;
(iii) strong and control variants within a family that differ only in numerical values, not template structure;
(iv) operands in control items that are ``hard numbers'' (last two digits in [25, 75], not divisible by 10, not near round boundaries), verified by automated checks.

The benchmark spans 4 digit scales ($d \in \{2, 4, 8, 16\}$), yielding $8 \text{ categories} \times 50 \text{ families} \times 4 \text{ scales} = 1{,}600$ families and $4{,}800$ individual items.
Scaling operand size increases computational load while leaving shortcut applicability unchanged in strong variants.

\subsection{Prompting Conditions}
\label{sec:conditions}

Each item is administered under three prompting conditions that differ in how strongly they encourage or suppress shortcut strategies:

\textbf{CoT} (chain-of-thought): standard step-by-step reasoning instruction.

\textbf{NS} (number-sense): encourages mathematical intuition and easy calculations. Crucially, the NS prompt does \emph{not} name any specific shortcut type; it is deliberately category-agnostic.

\textbf{Strict}: explicitly forbids shortcuts, requiring digit-by-digit computation. Serves as a negative control.

All conditions require answers inside \texttt{\textbackslash boxed\{\}} for unambiguous extraction. Full prompt templates are provided in Appendix~\ref{sec:appendix_prompts}.

\subsection{Evaluation Metrics}
\label{sec:metrics}

\textbf{Accuracy.} The primary metric is per-variant accuracy under each prompting condition. The key signature of number-sense exploitation is an asymmetric prompting effect: NS prompting maintains or improves accuracy on strong items while reducing accuracy on control items.

\textbf{NS rate.} To directly measure strategy choice beyond accuracy, we use GPT-4.1-mini as an automated judge to classify each model response as SHORTCUT or COMPUTATION. This provides a process-level measure of whether NS prompting actually changes the reasoning strategy, not just the final answer.

\textbf{Normalised improvement.} To compare across categories with different CoT baselines, we report the normalised improvement $(\text{NS} - \text{CoT}) / (1 - \text{CoT})$, which measures what fraction of the remaining accuracy gap NS captures. A value of $+1$ means NS closes the entire remaining gap; $0$ means no change; negative means NS hurts performance.

\subsection{Judge and Generate Tasks}
\label{sec:judge_gen}

Beyond Use-level accuracy, we evaluate declarative and constructive shortcut competence.

\textbf{Judge tasks.} J1 presents a problem with options and asks whether it could be solved faster with mental math (YES/NO). J2 presents a problem with a solution and asks whether the solution used a shortcut or computation strategy (SHORTCUT/COMPUTATION). These measure whether models can \emph{recognise} shortcut applicability (J1) and \emph{identify} strategy type (J2).

\textbf{Generate task (G2).} Given a category description and one example, models generate a new strong/control problem pair. Generated items are verified by 6 deterministic code checks (answer correctness, shortcut existence, control blocks shortcut, family matching, novelty, digit scale). This measures whether models can \emph{construct} shortcut-amenable problems.

Full task specifications and verification details are provided in Appendix~\ref{sec:appendix_tasks}.
